\section{Conclusion}
\label{sec:discussion}
\bench{} evaluates the connection between educational procedures, agent
behavior, and delivered teaching artifacts. Its 305 tasks and 14 Skill
families support paired outcome comparisons, observable package access, and
family-level diagnosis. The results show useful procedural detail in selected
cases alongside incomplete delivery and resource failures after access.
Small-model artifact learning provides a complementary route to studying
performance without external documents. Expert-adjudicated breakdowns clarify
where reference-conditioned judgments agree and where their components differ.
Together, these analyses turn a coarse question---whether Skills help---into
testable questions about access, appropriate operations, and usable products.
\label{page:main-end}

\section*{Limitations}
The 305-task inventory is unbalanced: lesson building contributes over half
the tasks, while five families contain three each. Task origin, language, and
rubric style are correlated. Access is self-selected, so loaded versus
non-loaded groups do not estimate a causal effect of reading a Skill. Full
raw traces are locally available for two native models; the other two have
retained tool summaries. Provider implementations and a single retained outcome
per cell limit conclusions about repeated-run variability.

Our fixed reference can contain errors, and joint scoring can depend on its
paired candidate. The 28-task review covers all Skills but is not a
proportionate sample of the 305-task inventory. Two experts jointly approved
one label set, providing no independent inter-rater reliability estimate.
Benchmark-derived training variants test related-family transfer rather than
unseen-family generalization; the latest small-model recovery aggregates still
require item-level reconciliation. Final-artifact supervision does not directly
teach native package calls. None of these scores measures student learning
outcomes or establishes deployment safety.

\section*{Ethical Considerations}
The study uses public procedural resources and de-identified teaching
situations. Generated artifacts can contain misleading examples or
inappropriate advice and require substantive review before classroom use.
Release materials retain source attribution and evaluation provenance while
omitting credentials and reviewer affiliations. The expert review procedure
and prior assistance in preparing its labels are documented in
Appendix~\ref{app:human}.
